% Mathématiques 6e — cours V3
% Source LaTeX rendue côté Astro/KaTeX.

\section{Objectifs}

Dans ce chapitre, on apprend à :

\begin{itemize}
\item choisir une unité de longueur adaptée ;
\item convertir des longueurs sans modifier la grandeur mesurée ;
\item mesurer et reporter une longueur avec une règle ou un compas ;
\item distinguer une longueur, un périmètre et une aire ;
\item calculer le périmètre de polygones usuels et de figures composées ;
\item comprendre le lien entre le diamètre d’un disque et son périmètre ;
\item utiliser la formule du périmètre d’un disque ;
\item résoudre des problèmes de longueur en contrôlant les unités et l’ordre de grandeur.
\end{itemize}

\section{1. Une longueur est une grandeur}

Une longueur permet de mesurer une distance, un segment, une hauteur, une largeur, un contour ou un trajet.

Une mesure de longueur comporte toujours :

\begin{itemize}
\item une \textbf{valeur numérique} ;
\item une \textbf{unité}.
\end{itemize}

Par exemple :

\[
2{,}4\ \text{m}
\]

et

\[
240\ \text{cm}
\]

représentent exactement la même longueur.

Changer d’unité ne change donc pas l’objet mesuré : seule l’écriture de la mesure change.

\section{2. Les unités usuelles de longueur}

Les principales unités sont organisées autour du mètre.

\coursefigure{/images/mathematiques/6eme/longueurs-unites.svg}{Schéma des unités de longueur du kilomètre au millimètre.}{Les unités successives sont le kilomètre, l’hectomètre, le décamètre, le mètre, le décimètre, le centimètre et le millimètre.}

Les relations fondamentales sont :

\[
1\ \text{km}=1\,000\ \text{m},
\]

\[
1\ \text{m}=10\ \text{dm}=100\ \text{cm}=1\,000\ \text{mm}.
\]

Et :

\[
1\ \text{cm}=10\ \text{mm}.
\]

\section{3. Convertir une longueur}

Convertir consiste à écrire une même longueur dans une autre unité.

\subsection{Exemple}

Convertissons :

\[
3{,}25\ \text{m}
\]

en centimètres.

Comme :

\[
1\ \text{m}=100\ \text{cm},
\]

on obtient :

\[
3{,}25\times100=325.
\]

Donc :

\[
\boxed{3{,}25\ \text{m}=325\ \text{cm}}.
\]

\subsection{Conversion inverse}

Convertissons :

\[
485\ \text{cm}
\]

en mètres.

Comme :

\[
100\ \text{cm}=1\ \text{m},
\]

on calcule :

\[
485\div100=4{,}85.
\]

Donc :

\[
\boxed{485\ \text{cm}=4{,}85\ \text{m}}.
\]

\section{4. Méthode de conversion}

Pour convertir une longueur :

\begin{enumerate}
\item identifier l’unité de départ ;
\item identifier l’unité d’arrivée ;
\item déterminer le facteur de conversion ;
\item effectuer la multiplication ou la division ;
\item vérifier que l’ordre de grandeur est cohérent.
\end{enumerate}

Par exemple, passer des mètres aux millimètres doit produire une valeur numérique plus grande, car le millimètre est une unité plus petite.

Ainsi :

\[
2\ \text{m}=2\,000\ \text{mm}.
\]

\section{5. Choisir l’unité adaptée}

Une unité est choisie en fonction de l’ordre de grandeur de l’objet.

On mesure généralement :

\begin{itemize}
\item l’épaisseur d’un petit objet en millimètres ;
\item la longueur d’un cahier en centimètres ;
\item la longueur d’une salle en mètres ;
\item une distance entre deux villes en kilomètres.
\end{itemize}

Une distance de route écrite :

\[
4\,850\ \text{m}
\]

peut être plus lisiblement exprimée :

\[
4{,}85\ \text{km}.
\]

Les deux écritures sont exactes, mais l’une est souvent plus adaptée au contexte.

\section{6. Mesurer avec un instrument}

Une mesure est toujours liée à l’instrument utilisé.

Une règle graduée au millimètre permet de lire une longueur au millimètre près.

Le compas peut servir à \textbf{reporter une longueur} : on règle l’écartement sur une longueur donnée puis on le transporte ailleurs sans modifier cet écartement.

Le compas ne sert donc pas uniquement à tracer des cercles.

\section{7. Reporter une longueur avec le compas}

Supposons qu’un segment ait pour longueur :

\[
AB.
\]

On ouvre le compas pour que ses pointes coïncident avec $A$ et $B$.

Sans modifier l’ouverture, on peut reporter exactement cette longueur à partir d’un autre point $C$.

On construit alors un point $D$ tel que :

\[
CD=AB.
\]

Cette propriété est utile dans de nombreuses constructions géométriques.

\section{8. Le périmètre}

Le périmètre d’une figure est la \textbf{longueur de son contour}.

Il s’exprime donc avec une unité de longueur :

\[
\text{mm},\quad \text{cm},\quad \text{m},\quad \text{km},\ldots
\]

Le périmètre ne doit pas être confondu avec l’aire, qui mesure une surface.

\section{9. Périmètre d’un rectangle}

Pour un rectangle de longueur $L$ et de largeur $\ell$ :

\[
P=L+\ell+L+\ell.
\]

On regroupe :

\[
\boxed{P=2(L+\ell)}.
\]

\subsection{Exemple développé 1}

Pour :

\[
L=8{,}5\ \text{cm}
\]

et :

\[
\ell=3\ \text{cm},
\]

on calcule :

\[
P=2(8{,}5+3).
\]

Donc :

\[
P=2\times11{,}5=23.
\]

Ainsi :

\[
\boxed{P=23\ \text{cm}}.
\]

\section{10. Périmètre d’un carré}

Si le côté d’un carré mesure $c$, ses quatre côtés ont la même longueur.

Donc :

\[
P=c+c+c+c.
\]

Ainsi :

\[
\boxed{P=4c}.
\]

Pour :

\[
c=6{,}2\ \text{cm},
\]

on obtient :

\[
P=4\times6{,}2=24{,}8\ \text{cm}.
\]

\section{11. Périmètre d’un polygone quelconque}

Pour un polygone, on additionne les longueurs de tous les côtés appartenant au contour.

Par exemple, un triangle de côtés :

\[
4{,}5\ \text{cm},\qquad5\ \text{cm},\qquad6{,}2\ \text{cm}
\]

a pour périmètre :

\[
P=4{,}5+5+6{,}2.
\]

Donc :

\[
\boxed{P=15{,}7\ \text{cm}}.
\]

\section{12. Figures composées}

Une figure composée rassemble plusieurs formes simples.

Pour calculer son périmètre, il ne faut pas additionner toutes les longueurs visibles dans le dessin.

Il faut uniquement suivre \textbf{le contour extérieur}.

Une arête intérieure qui sépare deux parties n’appartient pas au périmètre extérieur.

\subsection{Méthode}

\begin{enumerate}
\item parcourir mentalement le contour ;
\item repérer uniquement les segments réellement extérieurs ;
\item déterminer les longueurs manquantes si nécessaire ;
\item convertir toutes les longueurs dans une même unité ;
\item additionner.
\end{enumerate}

\section{13. Exemple développé 2 — figure en forme de marche}

Une figure composée possède un contour constitué des longueurs :

\[
6\ \text{cm},\quad
2\ \text{cm},\quad
3\ \text{cm},\quad
4\ \text{cm},\quad
3\ \text{cm},\quad
6\ \text{cm}.
\]

Son périmètre vaut :

\[
P=6+2+3+4+3+6.
\]

Donc :

\[
\boxed{P=24\ \text{cm}}.
\]

La difficulté principale est souvent de sélectionner les bons segments, pas d’effectuer l’addition.

\section{14. Cercle et disque}

Un \textbf{cercle} est le contour.

Un \textbf{disque} est la surface intérieure délimitée par ce cercle.

Le périmètre du disque est donc la longueur du cercle qui le borde.

\coursefigure{/images/mathematiques/6eme/perimetre-disque-diametre.svg}{Disque montrant le centre, un rayon, un diamètre et le contour dont on calcule le périmètre.}{Le diamètre traverse le centre et vaut deux fois le rayon : $d=2r$.}

\section{15. Rayon et diamètre}

Si $r$ est le rayon et $d$ le diamètre :

\[
\boxed{d=2r}.
\]

Inversement :

\[
\boxed{r=\dfrac{d}{2}}.
\]

Il faut identifier correctement la donnée fournie avant d’utiliser une formule.

\section{16. Une proportionnalité remarquable}

Lorsque l’on mesure plusieurs disques, on constate que le quotient :

\[
\dfrac{\text{périmètre}}{\text{diamètre}}
\]

est toujours proche du même nombre.

Ce nombre est noté :

\[
\pi.
\]

On a approximativement :

\[
\pi\approx3{,}14.
\]

Le périmètre d’un disque est donc \textbf{proportionnel à son diamètre}.

\section{17. Formule du périmètre d’un disque}

Si $d$ est le diamètre :

\[
\boxed{P=\pi d}.
\]

Comme :

\[
d=2r,
\]

on peut aussi écrire :

\[
\boxed{P=2\pi r}.
\]

Les deux formules sont équivalentes.

\section{18. Exemple développé 3 — disque donné par son rayon}

Un disque a pour rayon :

\[
r=4\ \text{cm}.
\]

Son périmètre exact est :

\[
P=2\pi\times4.
\]

Donc :

\[
P=8\pi\ \text{cm}.
\]

Avec :

\[
\pi\approx3{,}14,
\]

on obtient :

\[
P\approx8\times3{,}14=25{,}12\ \text{cm}.
\]

Ainsi :

\[
\boxed{P\approx25{,}1\ \text{cm}}.
\]

\section{19. Exemple développé 4 — disque donné par son diamètre}

Un disque a pour diamètre :

\[
d=15\ \text{cm}.
\]

On utilise directement :

\[
P=\pi d.
\]

Donc :

\[
P=15\pi\ \text{cm}.
\]

Une valeur approchée est :

\[
P\approx15\times3{,}14=47{,}1\ \text{cm}.
\]

\section{20. Contrôler un périmètre de disque}

Le périmètre doit être un peu plus de trois fois le diamètre.

Si :

\[
d=10\ \text{cm},
\]

le périmètre doit être proche de :

\[
3\times10=30\ \text{cm}.
\]

Un résultat de :

\[
3{,}1\ \text{cm}
\]

ou de :

\[
310\ \text{cm}
\]

serait donc immédiatement suspect.

\section{21. Mélanger des unités}

Avant d’additionner des longueurs, elles doivent être exprimées dans une même unité.

Supposons un rectangle de dimensions :

\[
1{,}2\ \text{m}
\]

et :

\[
75\ \text{cm}.
\]

On convertit :

\[
75\ \text{cm}=0{,}75\ \text{m}.
\]

Puis :

\[
P=2(1{,}2+0{,}75).
\]

Donc :

\[
P=2\times1{,}95=3{,}9\ \text{m}.
\]

\section{22. Résoudre un problème de clôture}

Un jardin rectangulaire mesure :

\[
18\ \text{m}
\]

sur :

\[
12\ \text{m}.
\]

Son périmètre total vaut :

\[
P=2(18+12)=60\ \text{m}.
\]

Si une ouverture de :

\[
3\ \text{m}
\]

ne doit pas être clôturée, la longueur nécessaire vaut :

\[
60-3=57\ \text{m}.
\]

Donc :

\[
\boxed{57\ \text{m}}
\]

de clôture sont nécessaires.

\section{23. Méthode — résoudre un problème de longueur}

\subsection{Étape 1 — identifier la grandeur cherchée}

S’agit-il :

\begin{itemize}
\item d’une longueur ;
\item d’un périmètre ;
\item d’un rayon ;
\item d’un diamètre ?
\end{itemize}

\subsection{Étape 2 — relever les données utiles}

On peut annoter la figure.

\subsection{Étape 3 — harmoniser les unités}

Aucune addition de longueurs dans des unités différentes.

\subsection{Étape 4 — choisir le calcul}

Formule, addition du contour, proportionnalité ou conversion.

\subsection{Étape 5 — contrôler}

Comparer le résultat avec un ordre de grandeur réaliste.

\subsection{Étape 6 — rédiger}

Donner la valeur avec l’unité adaptée.

\section{24. Erreurs fréquentes}

\subsection{Confondre périmètre et aire}

Le périmètre mesure un contour.

Il s’exprime en unité de longueur, par exemple :

\[
\text{cm}.
\]

Une aire s’exprime en unité carrée, par exemple :

\[
\text{cm}^2.
\]

\subsection{Additionner des unités différentes}

Une écriture comme :

\[
2\ \text{m}+35\ \text{cm}
\]

doit être convertie avant de donner une somme numérique unique.

\subsection{Confondre rayon et diamètre}

Si la donnée est le diamètre, ne pas le multiplier encore par :

\[
2.
\]

\subsection{Utiliser la formule du disque pour un demi-cercle sans réfléchir}

Dans une figure composée, il faut distinguer la longueur de l’arc et les segments droits présents dans le contour.

\subsection{Oublier une ouverture}

Dans un problème de clôture ou de bordure, il faut lire précisément quelles parties du contour sont réellement concernées.

\section{25. À retenir}

Une conversion change l’unité, pas la longueur.

Le périmètre est la longueur du contour.

Pour un rectangle :

\[
\boxed{P=2(L+\ell)}.
\]

Pour un carré :

\[
\boxed{P=4c}.
\]

Pour un disque :

\[
\boxed{P=\pi d=2\pi r}.
\]

Le périmètre d’un disque est proportionnel à son diamètre.

Avant tout calcul de périmètre :

\begin{itemize}
\item mettre les longueurs dans la même unité ;
\item identifier le contour réel ;
\item choisir la formule ou l’addition adaptée ;
\item vérifier l’ordre de grandeur ;
\item donner une unité de longueur.
\end{itemize}
